\chapter*{Glossary}
\addcontentsline{toc}{chapter}{Glossary}

\begin{description}[style=nextline,leftmargin=1.6em,itemsep=0.75em]
  \item[Absolute path] A file location beginning from a file-system root or
        user-specific location. It is normally unsuitable for a portable
        project.
  \item[Accessibility] The practice of making content usable by people with
        varied sensory, motor, and cognitive access needs, including structured
        navigation and meaningful alternatives to visual information.
  \item[Affiliation] The institution or organisation associated with an author
        under the target venue's rules.
  \item[Argument] Information supplied to a command. Required arguments are
        commonly placed in braces; optional arguments in square brackets.
  \item[Auxiliary file] A generated file used between compilation passes to
        store references, contents, citations, index entries, or other state.
  \item[Backend] A separate program used by a \LaTeX{} package, such as Biber
        processing bibliography control data for \pkg{biblatex}.
  \item[Back matter] Material after the main body, such as references, glossary,
        solutions, index, or author information.
  \item[Baseline] The imaginary line on which most letters sit; also a known
        version used as a comparison point in testing.
  \item[Bibliography database] A structured text file, commonly ending in
        \file{.bib}, that stores verified source metadata under citation keys.
  \item[Bookmark] A navigational entry in a PDF outline, usually derived from a
        structural heading.
  \item[Build] The complete process that transforms source and dependencies into
        output, including repeat \LaTeX{}, bibliography, glossary, or index passes.
  \item[Build log] The engine's technical record of files, messages, warnings,
        errors, fonts, and output from a compilation.
  \item[Caption] Text that identifies and explains a figure or table sufficiently
        for its role in the document.
  \item[Checksum] A computed digest used to detect whether a file's bytes have
        changed; it does not establish scientific validity or authorship.
  \item[Citation] A reference marker in the text connecting a claim or discussion
        to a bibliography record.
  \item[Citation key] The unique database identifier used inside a citation
        command, for example \code{lamport1994}.
  \item[Class option] A supported choice passed in square brackets to
        \cmd{documentclass}, such as base size or two-sided layout.
  \item[Clean build] Compilation after removing or isolating generated state so
        that missing dependencies and stale outputs cannot hide easily.
  \item[Command] A \LaTeX{} instruction, normally beginning with a backslash.
  \item[Compilation] Interpretation of source to produce output and diagnostic
        records.
  \item[Compiler or engine] The program that reads TeX/\LaTeX{} source, such as
        pdfTeX, XeTeX, or LuaTeX in an appropriate format.
  \item[Compliance matrix] A table mapping each external requirement to its
        implementation and verification evidence.
  \item[Conditional] Source logic that selects one branch, such as an anonymous
        or identified author block.
  \item[Continuous integration] An automated service that builds or tests a
        versioned project after defined events.
  \item[Controlled vocabulary] An authorised set of terms used for consistent
        indexing, metadata, or domain description.
  \item[Cross-reference] A generated connection to a numbered or named target
        such as a figure, equation, section, or appendix.
  \item[Data dictionary] A structured explanation of variables, units, codes,
        formats, missing values, and constraints in a data set.
  \item[Declaration] A formal statement about matters such as funding, interests,
        ethics, authorship, originality, or data access. Required wording must be
        verified.
  \item[Dependency] A class, package, font, program, asset, or other object
        required to reproduce a build.
  \item[Document body] Content between
        \cmd{begin}\required{document} and \cmd{end}\required{document}.
  \item[Document class] The broad structural and visual foundation selected by
        \cmd{documentclass}.
  \item[Document preparation system] A system in which authors describe content
        and structure from which formatted output is constructed.
  \item[Effective resolution] Raster pixels divided by the image's final physical
        dimensions, commonly expressed as pixels per inch.
  \item[Engine] See \emph{compiler or engine}.
  \item[Environment] A structured region opened by \cmd{begin} and closed by a
        matching \cmd{end}.
  \item[Error] A condition that prevents the engine from interpreting source
        reliably; the first meaningful error is usually the best starting point.
  \item[Evidence register] A controlled list of data, code, figures, reviews,
        approvals, and other objects supporting a report or release.
  \item[Float] A figure or table object whose final placement \LaTeX{} chooses under
        declared preferences and page constraints.
  \item[Font embedding] Inclusion of a font or subset inside the PDF so the
        viewer need not obtain it separately. Embedding does not settle licence.
  \item[Front matter] Preliminary material such as title, declarations, abstract,
        contents, and lists, often with distinct numbering.
  \item[Generated file] An output created from source, such as a PDF figure,
        table fragment, auxiliary record, or compiled document.
  \item[Glossary] An alphabetical or structured list defining specialised terms
        used in a document.
  \item[Identifier] A stable name assigned to an object, such as a DOI, accession
        number, report ID, figure label, or repository record.
  \item[Index] An editorial map from reader search terms to useful pages, not
        merely an alphabetical word list.
  \item[Input file] A source, data, bibliography, figure, or configuration file
        consumed by a build.
  \item[Interoperability] The capacity of source and outputs to work across
        relevant systems, tools, or repositories under documented conditions.
  \item[Label] A unique source key attached to a numbered target for later
        cross-reference.
  \item[\LaTeX{} project] The complete group of source, bibliography, figure, data,
        configuration, documentation, permission, and build files needed for a
        defined output.
  \item[Licence] Terms granting specified permissions to use, modify, or share a
        copyrighted work. A public URL is not itself a licence.
  \item[Line breaking] The typographic process of distributing words across
        lines under width, hyphenation, and spacing constraints.
  \item[Log] See \emph{build log}.
  \item[Main document] The source file given to the engine first; it normally
        owns the class, shared setup, and document order.
  \item[Manifest] A list of files in an exact release, sometimes accompanied by
        sizes, versions, checksums, origins, or access classes.
  \item[Metadata] Structured descriptive information such as title, authors,
        abstract, keywords, version, date, identifiers, and rights.
  \item[Minimal working example (MWE)] The smallest complete source that still
        reproduces a behaviour or problem.
  \item[Monospaced] A font design in which characters normally occupy equal
        horizontal width, often used to distinguish source and filenames.
  \item[Nomenclature] An organised list of symbols and their meanings.
  \item[Option] A command or package choice commonly placed in square brackets;
        accepted options are defined by the receiving interface.
  \item[Orphan] A typographic term often used for the first line of a paragraph
        left at the bottom of a page; definitions vary across style guides.
  \item[Output directory] The location assigned to generated build files.
  \item[Overfull box] A TeX box whose content exceeds its allocated dimension,
        potentially producing text beyond a margin.
  \item[Package] A specialised extension loaded for capabilities not supplied by
        the selected class and \LaTeX{} kernel alone.
  \item[Page anchor] An internal PDF destination associated with a page. Duplicate
        page numbers need distinct anchors.
  \item[Persistent identifier] A managed identifier intended to remain stable,
        such as a DOI; persistence depends on its governance and metadata.
  \item[Placeholder] An explicit marker for information that is unavailable or
        unverified. It must be resolved or deliberately retained before release.
  \item[Preamble] Source before \cmd{begin}\required{document}, containing the
        class and document-wide setup.
  \item[Proof] A document version intended for checking rather than final
        distribution; a submission-system proof is a newly generated output.
  \item[Provenance] The documented origin, transformations, responsibility, and
        history of data, code, figures, or claims.
  \item[Raster image] An image represented by pixels; useful for photographs and
        microscopy but resolution-dependent at placed size.
  \item[Reference list] The formatted collection of cited source records.
  \item[Relative path] A file location expressed from the project or current
        document directory, enabling portability when the structure is preserved.
  \item[Release] An identified, controlled set of source and outputs intended for
        a recipient, submission, deposit, or publication.
  \item[Repository] A managed storage and discovery location for source, data,
        software, documents, or records, with its own policies and identifiers.
  \item[Reproducible compilation] The ability of an authorised reader to obtain
        the stated document output from preserved source, dependencies, and
        instructions.
  \item[Reserved character] A character with a syntactic role in \LaTeX{} source
        that needs special treatment when printed literally.
  \item[Resolution] The density or number of samples in an image. Venue rules
        usually distinguish raster type and final placed size.
  \item[Semantic command] A custom command named for meaning, allowing notation
        or presentation to be changed consistently in one definition.
  \item[Shell escape] A TeX capability that permits external commands. It can
        support some packages but expands the build's security surface.
  \item[Source file] Editable plain text containing content and \LaTeX{}
        instructions, commonly with extension \file{.tex}.
  \item[Structured data] Data organised under a documented schema so variables
        and records can be interpreted and processed consistently.
  \item[Submission package] The exact set of files sent to a journal,
        institution, repository, or client under its current instructions.
  \item[Supplementary information] Material distributed alongside a main work,
        governed by the recipient's format, numbering, review, and access rules.
  \item[TeX distribution] A maintained collection of engines, formats, packages,
        fonts, documentation, and utilities.
  \item[Toolchain] The ordered set of programs and configurations that transforms
        inputs into outputs.
  \item[Underfull box] A box whose content cannot be spaced to the desired
        dimension without looseness; it requires visual judgement.
  \item[Vector graphic] An image represented by paths and objects, normally
        scalable for line art, diagrams, plots, and text.
  \item[Version control] A system that records identified snapshots and changes,
        allowing comparison, collaboration, and recovery.
  \item[Warning] A diagnostic condition that permits output but requires
        judgement, additional passes, or visual inspection.
  \item[Widow] A typographic term often used for the final line of a paragraph
        isolated at the top of a page; definitions vary across style guides.
  \item[Working directory] The directory from which a process resolves relative
        paths. The project should document its expected main build directory.
\end{description}
